\documentclass[a4paper,12pt,titlepage,finall]{article}

\usepackage[T1,T2A]{fontenc}     % форматы шрифтов
\usepackage[utf8x]{inputenc}     % кодировка символов, используемая в данном файле
\usepackage[russian]{babel}      % пакет русификации
\usepackage{tikz}                % для создания иллюстраций
\usepackage{pgfplots}            % для вывода графиков функций
\usepackage{geometry}		 % для настройки размера полей
\usepackage{indentfirst}         % для отступа в первом абзаце секции
\usepackage{minted}
\usepackage{amsmath}
\usepackage{url}

% выбираем размер листа А4, все поля ставим по 3см
\geometry{a4paper,left=30mm,top=30mm,bottom=30mm,right=30mm}

\setcounter{secnumdepth}{0}      % отключаем нумерацию секций

\usepgfplotslibrary{fillbetween} % для изображения областей на графиках

\begin{document}
% Титульный лист
\begin{titlepage}
    \begin{center}
	{\small \sc Московский государственный университет \\имени М.~В.~Ломоносова\\
	Факультет вычислительной математики и кибернетики\\}
	\vfill
	{\Large \sc Отчет по заданию №6}\\
	~\\
	{\large \bf <<Сборка многомодульных программ. \\
	Вычисление корней уравнений и определенных интегралов.>>}\\ 
	~\\
	{\large \bf Вариант 11 / 5 / 4}
    \end{center}
    \begin{flushright}
	\vfill {Выполнил:\\
	студент 102 группы\\
	Соловьев~А.~И.\\
	~\\
	Преподаватели:\\
	Калинин~Г.~М.\\Цыбров~Е.~Г.}
    \end{flushright}
    \begin{center}
	\vfill
	{\small Москва\\2024}
    \end{center}
\end{titlepage}

% Автоматически генерируем оглавление на отдельной странице
\tableofcontents
\newpage

\section{Постановка задачи}

Необходимо найти площадь фигуры, ограниченной данными тремя кривыми, лежащей на заданном отрезке. Гарантируется, что кривые попарно пересекаются на данном отрезке. 
\begin{itemize}
    \item Для нахождения точек пересечений требуется реализовать численный метод нахождения корней. Решение реализует два метода (на выбор): метод Ньютона (касательных) и метод хорд.
    \item Для нахождения площади фигуры требуется реализовать метод численного интегрирования. Решение реализует два метода (на выбор): метод трапеций и метод Симпсона (парабол).
\item Отрезок для применения метода нахождения корней должен быть вычислен аналитически.
\end{itemize}
Функции, задающие кривые, и отрезок, на котором лежит фигура, задаются перед сборкой в файле. Функции задаются в обратной польской записи.\\\\
Пример:
\begin{minted}[frame=single, linenos, breaklines=true, fontsize=\small]{text}
0.0 4.0
2 x 4 / tan -
x
0.2 pi *
\end{minted}

\newpage

\section{Математическое обоснование}

\subsection{Минимальные условия сходимости}
\begin{itemize}
    \item $\exists x_0: f(x_0) = 0$
    \item $f(x)$ непрерывно дифференцируема на отрезке
    \item Начальная точка $x$ достаточна близка к $x_0$
\end{itemize}

Для метода Ньютона желательно соблюсти все условия Фурье:
\begin{itemize}
    \item $f'(x) \neq 0$ на отрезке
    \item $f''(x)$ не меняет знак на отрезке
    \item начальная точка $x_0$ с того конца отрезка, где $f(x_0)\cdot f''(x_0) > 0$,
\end{itemize}
так как при соблюдении лишь минимальных условий метод может уйти за пределы отрезка, сойтись к другому корню или не сойтись вовсе.~\cite{math}\\\\

Площадь заданной фигуры должна быть вычислена с точность $\varepsilon$. В некоторой окрестности точки пересечения, которую мы ищем с точностью до $\varepsilon_1$, функции эквивалентны своим касательным в этой точке.~\cite{math} Значит, при смещении границы интегрирования, изменение значения интеграла будет пропорционально квадрату смещения.
Общая площадь фигуры равна комбинации трех интегралов с некоторыми знаками; общая погрешность может достигать $3\varepsilon_2 + 2\varepsilon_1^2 \leq 0.001$ по условию. Возьмем $\varepsilon_2 = \varepsilon_1 = 0.0001$, $3\varepsilon_2 + 2\varepsilon_1^2 = 0.0003 + 0.0000002 \approx 0.0003 \leq 0.001$.\\

Для вычисления интеграла с определенной точностью было использовано правило Рунге. Метод Симпсона имеет порядок точности 4, метод трапеций имеет порядок точности 15. Оценка погрешности равна разнице между двумя площадями, посчитанными с разной точностью, деленной на порядок точности. Поэтому количество разбиений в обоих методах экспоненциально повышается, пока оценка не станет меньше $\varepsilon_2$.\\

После нахождения точек, сначала определяется верхняя/нижняя ограничивающая кривая. Это кривая, которая не пересекает другие в средней (по оси $x$) точке пересечения. Определив, лежит ли она ниже или выше этой точки, мы либо вычитаем ее интеграл из интегралов двух других функций, либо наоборот:\\\\
$S = sign \cdot (\int_a^b f_1(x)\,dx - \int_a^{x_1} f_2(x)\,dx - \int_{x_1}^b f_3(x)\,dx)$,\\\\
где sign равен 1, если кривая выше $x_1$, и -1, если ниже. 

\newpage

\section{Результаты тестов}

Программа тестировалась на таких наборах данных:\\
1.
\begin{minted}[frame=single, linenos, breaklines=true, fontsize=\small]{text}
0.0 1.0
x
1 x -
x 0.5 - 2 ^
\end{minted}
Аналитическое решение:
$$
\begin{gathered}
    x^2 - x + 0.25 = x\\D = (-2)^2-4\cdot0.25=3\\x_{1,2}=\frac{2\pm \sqrt{3}}{2}=1 \pm \frac{\sqrt{3}}{2},\; x_1 \in [0,\,1]\\\\
    x^2 - x + 0.25 = 1-x,\,x^2-0.75=0\\D = 4\cdot0.75=3\\x_{1,2}=\frac{\pm \sqrt{3}}{2},\;x_2 \in [0, 1]\\\\
x = 1-x;\; x_3=0.5
\end{gathered}
$$

$$
\begin{gathered}
\int x\,dx = \frac{1}{2} x^2 + C\\
\int -x+1\,dx = -\frac{1}{2}x^2+x + C\\
\int (x-0.5)^2\,dx = \frac{1}{3}(x-0.5)^3 + C\\\\
\end{gathered}
$$

$$
\begin{gathered}
\frac{1}{2}\cdot(0.25 - (1 - \frac{\sqrt{3}}{2})^2) \approx 0.116\\
-\frac{1}{2}\cdot(\frac{3}{4} - 0.25) + (\frac{\sqrt{3}}{2} - 0.5) \approx 0.366 - 0.25 = 0.116\\
\frac{1}{3}\cdot (0.5 - \frac{\sqrt{3}}{2})^3 \approx 0.016\\
\frac{1}{3}\cdot (\frac{\sqrt{3}}{2}-0.5)^3 \approx 0.016\\\\
S \approx 0.116 \cdot 2 - 0.016 \cdot 2 = 0.2
\end{gathered}
$$
\\\\


\newpage


2.
\begin{minted}[frame=single, linenos, breaklines=true, fontsize=\small]{text}
0.0 1.0
x tan
1 x tan -
x tan 10 /
\end{minted}
Аналитическое решение:
$$
\begin{gathered}
tan(x) = 1 - tan(x);\;tan(x) = 0.5;\; x_1 = arctan(0.5);\\
tan(x) = \frac{tan(x)}{10};\;tan(x) = 0; x_2 = 0;\\
1 - tan(x) = \frac{tan(x)}{10};\;tan(x)=\frac{10}{11}; x_3 = arctan(\frac{10}{11})\\
\end{gathered}
$$

$$
\begin{gathered}
\int tan(x)\,dx = -ln|cos\,x| + C\\
\int 1 - tan(x)\,dx = x + ln|cos\;x|+C\\
\int \frac{tan(x)}{10}\,dx = \frac{-ln|cos\,x|}{10} + C\\\\
S = -ln|cos(arctan(0.5)| + (arctan(\frac{10}{11}) + ln|cos(arctan(\frac{10}{11})| - \\- (arctan(0.5)+ln|cos(arctan(0.5|)) + \frac{ln|cos(arctan(\frac{10}{11}))|}{10} = \\
\approx 0.166
\end{gathered}
$$


\newpage
Результаты вычислений для двух вышеупомянутых тестовых примеров:
координаты точек пересечения в таблицах и площадь полученной фигуры.

\begin{table}[h]
\centering
\begin{tabular}{|c|c|c|}
\hline
Кривые & $x$ & $y$ \\
\hline
1 и 2 &  0.5000 & 0.5000 \\
2 и 3 &  0.8660 & 0.1339 \\
1 и 3 & 0.1339 & 0.1339 \\
\hline
\end{tabular}
\caption{Координаты точек пересечения}
\label{table1}
\end{table}


\begin{figure}[h]
\centering
\begin{tikzpicture}
\begin{axis}[% grid=both,                % рисуем координатную сетку (если нужно)
             axis lines=middle,          % рисуем оси координат в привычном для математики месте
             restrict x to domain=-1:2,  % задаем диапазон значений переменной x
             restrict y to domain=-0.25:1.5,  % задаем диапазон значений функции y(x)
             axis equal,                 % требуем соблюдения пропорций по осям x и y
             enlargelimits,              % разрешаем при необходимости увеличивать диапазоны переменных
             legend cell align=left,     % задаем выравнивание в рамке обозначений
             scale=2,                    % задаем масштаб 2:1
             xticklabels={,,},           % убираем нумерацию с оси x
             yticklabels={,,}]           % убираем нумерацию с оси y

% первая функция
% параметр samples отвечает за качество прорисовки
\addplot[green,samples=256,thick,name path=A] {x};
% описание первой функции
    \addlegendentry{$f_1(x)=x$}

% добавим немного пустого места между описанием первой и второй функций
\addlegendimage{empty legend}\addlegendentry{}

% вторая функция
% здесь необходимо дополнительно ограничить диапазон значений переменной x
\addplot[blue,samples=256,thick,name path=B] {1-x};
    \addlegendentry{$f_2(x)=1-x$}

% дополнительное пустое место не требуется, так как формулы имеют небольшой размер по высоте

% третья функция
    \addplot[red,samples=256,thick,name path=C] {(x-0.5)^2};
    \addlegendentry{$f_3(x)=(x-0.5)^2$}

% закрашиваем фигуру
\addplot[blue!20,samples=256] fill between[of=B and C,soft clip={domain=0.5:0.8660}];
\addplot[blue!20,samples=256] fill between[of=A and C,soft clip={domain=0.1339:0.5}];
\addlegendentry{$S=0.1994$}

% Поскольку автоматическое вычисление точек пересечения кривых в TiKZ реализовать сложно,
% будем явно задавать координаты.
\addplot[dashed] coordinates { (0.1339, 0.1339) (0.1339, 0) };
\addplot[color=black] coordinates {(0.1339, 0)} node [label={-135:{\small 0.1339}}]{};

\addplot[dashed] coordinates { (0.5, 0.5) (0.5, 0) };
\addplot[color=black] coordinates {(0.5, 0)} node [label={-90:{\small 0.5}}]{};

\addplot[dashed] coordinates { (0.8660, 0.1339) (0.8660, 0) };
\addplot[color=black] coordinates {(0.8660, 0)} node [label={-45:{\small 0.8660}}]{};

\end{axis}
\end{tikzpicture}
\caption{Плоская фигура, ограниченная графиками заданных уравнений}
\label{plot2}
\end{figure}




















\begin{table}[h]
\centering
\begin{tabular}{|c|c|c|}
\hline
Кривые & $x$ & $y$ \\
\hline
1 и 2 &  0.4636 & 0.5000 \\
2 и 3 &  0.7378 & 0.091 \\
1 и 3 & 0.0000 & 0.0000 \\
\hline
\end{tabular}
\caption{Координаты точек пересечения}
\label{table1}
\end{table}

\begin{figure}[h]
\centering
\begin{tikzpicture}
\begin{axis}[% grid=both,                % рисуем координатную сетку (если нужно)
             axis lines=middle,          % рисуем оси координат в привычном для математики месте
             restrict x to domain=-1:2,  % задаем диапазон значений переменной x
             restrict y to domain=-0.25:1.5,  % задаем диапазон значений функции y(x)
             axis equal,                 % требуем соблюдения пропорций по осям x и y
             enlargelimits,              % разрешаем при необходимости увеличивать диапазоны переменных
             legend cell align=left,     % задаем выравнивание в рамке обозначений
             scale=2,                    % задаем масштаб 2:1
             xticklabels={,,},           % убираем нумерацию с оси x
             yticklabels={,,}]           % убираем нумерацию с оси y

% первая функция
% параметр samples отвечает за качество прорисовки
    \addplot[green,samples=256,thick,name path=A] {tan(x * 180 / 3.14159)};
% описание первой функции
    \addlegendentry{$f_1(x)=tan(x)$}

% добавим немного пустого места между описанием первой и второй функций
\addlegendimage{empty legend}\addlegendentry{}

% вторая функция
% здесь необходимо дополнительно ограничить диапазон значений переменной x
    \addplot[blue,samples=256,thick,name path=B] {1-tan(x * 180 / 3.14159)};
    \addlegendentry{$f_2(x)=1-tan(x)$}

% дополнительное пустое место не требуется, так как формулы имеют небольшой размер по высоте

% третья функция
    \addplot[red,samples=256,thick,name path=C] {tan(x * 180 / 3.14159)/10};
    \addlegendentry{$f_3(x)=\frac{tan(x)}{10}$}

% закрашиваем фигуру
\addplot[blue!20,samples=256] fill between[of=A and C,soft clip={domain=0.0:0.4636}];
\addplot[blue!20,samples=256] fill between[of=B and C,soft clip={domain=0.4636:0.7378}];
\addlegendentry{$S=0.166$}

% Поскольку автоматическое вычисление точек пересечения кривых в TiKZ реализовать сложно,
% будем явно задавать координаты.
\addplot[color=black] coordinates {(0.0, 0)} node [label={-135:{\small 0.0}}]{};

\addplot[dashed] coordinates { (0.4636, 0.5000) (0.4636, 0) };
\addplot[color=black] coordinates {(0.4636, 0)} node [label={-90:{\small 0.4636}}]{};

\addplot[dashed] coordinates { (0.7378, 0.091) (0.7378, 0) };
\addplot[color=black] coordinates {(0.7378, 0)} node [label={-45:{\small 0.7378}}]{};

\end{axis}
\end{tikzpicture}
\caption{Плоская фигура, ограниченная графиками заданных уравнений}
\label{plot2}
\end{figure}















\newpage
Также, помимо аналитически решенных систем, были протестированы те, которые аналитически не решаются. Например:
\begin{minted}[frame=single, linenos, breaklines=true, fontsize=\small]{text}
0.0 4.0
2 x 4 / tan -
x
0.6 pi *
\end{minted}
и
\begin{minted}[frame=single, linenos, breaklines=true, fontsize=\small]{text}
0.0 4.0
2 x 4 / tan -
x
0.6 pi *
\end{minted}
\begin{table}[h]
\centering
\begin{tabular}{|c|c|c|}
\hline
Кривые & $x$ & $y$ \\
\hline
1 и 2 &  1.5823 & 1.5823 \\
2 и $3_1$ &  2.557 & 1.256 \\
1 и $3_1$ & 1.256 & 1.256 \\
2 и $3_2$ &  0.4581 & 1.885 \\
1 и $3_2$ & 1.885 & 1.885 \\
\hline
\end{tabular}
\caption{Координаты точек пересечения}
\label{table1}
\end{table}

\begin{figure}[h]
\centering
\begin{tikzpicture}
\begin{axis}[% grid=both,                % рисуем координатную сетку (если нужно)
             axis lines=middle,          % рисуем оси координат в привычном для математики месте
             restrict x to domain=-1:3,  % задаем диапазон значений переменной x
             restrict y to domain=-0.25:4.5,  % задаем диапазон значений функции y(x)
             axis equal,                 % требуем соблюдения пропорций по осям x и y
             enlargelimits,              % разрешаем при необходимости увеличивать диапазоны переменных
             legend cell align=left,     % задаем выравнивание в рамке обозначений
             scale=1.9,                    % задаем масштаб 2:1
             xticklabels={,,},           % убираем нумерацию с оси x
             yticklabels={,,}]           % убираем нумерацию с оси y

% первая функция
% параметр samples отвечает за качество прорисовки
    \addplot[green,samples=256,thick,name path=A] {x};
% описание первой функции
    \addlegendentry{$f_1(x)=x$}

% добавим немного пустого места между описанием первой и второй функций
\addlegendimage{empty legend}\addlegendentry{}

% вторая функция
% здесь необходимо дополнительно ограничить диапазон значений переменной x
    \addplot[blue,samples=256,thick,name path=B] {2-tan((x * 180 / 3.14159) / 4)};
    \addlegendentry{$f_2(x)=2-tan(\frac{x}{4})$}

% дополнительное пустое место не требуется, так как формулы имеют небольшой размер по высоте

% третья функция
    \addplot[red,samples=256,thick,name path=C] {0.4 * 3.14159};
    \addlegendentry{$f_{3_1}(x)=0.4 \cdot \pi$}

% добавим немного пустого места между описанием третьей и четвертой функций
\addlegendimage{empty legend}\addlegendentry{}

% четвертая функция
    \addplot[yellow,samples=256,thick,name path=D] {0.6 * 3.14159};
    \addlegendentry{$f_{3_2}(x)=0.6 \cdot \pi$}

% закрашиваем фигуру
\addplot[blue!20,samples=256] fill between[of=A and C,soft clip={domain=1.256:1.5823}];
\addplot[blue!20,samples=256] fill between[of=B and C,soft clip={domain=1.5823:2.557}];
\addlegendentry{$\textcolor{blue}{S_1} = \textcolor{red}{S_2} = 0.166$}
В данном разделе приведены р
% закрашиваем фигуру
\addplot[red!20,samples=256] fill between[of=B and D,soft clip={domain=0.4581:1.5823}];
\addplot[red!20,samples=256] fill between[of=A and D,soft clip={domain=1.5823:1.885}];

% Поскольку автоматическое вычисление точек пересечения кривых в TiKZ реализовать сложно,
% будем явно задавать координаты.
\addplot[color=black] coordinates {(0.0, 0)} node [label={-135:{\small 0.0}}]{};

\addplot[dashed] coordinates { (0.4581, 1.885) (0.4581, 0) };
\addplot[color=black] coordinates {(0.4581, 0)} node [label={-90:{\small 0.4581}}]{};

\addplot[dashed] coordinates { (1.256, 1.256) (1.256, 0) };
\addplot[color=black] coordinates {(1.256, 0)} node [label={-90:{\small 1.256}}]{};


\addplot[dashed] coordinates { (2.557, 1.256) (2.557, 0) };
\addplot[color=black] coordinates {(2.557, 0)} node [label={-45:{\small 2.557}}]{};

\addplot[dashed] coordinates { (1.885, 1.885) (1.885, 0) };
\addplot[color=black] coordinates {(1.885, 0)} node [label={-90:{\small 1.885}}]{};

\end{axis}
\end{tikzpicture}
\caption{Плоская фигура, ограниченная графиками заданных уравнений}
\label{plot2}
\end{figure}

Этот тест показывает, что площадь вычисляется одинаково независимо от положения фигуры в пространстве. 


\newpage

\section{Структура программы и спецификация функций}

Для ознакомления с кодом см. репозиторий [3] и документацию [2].

\subsection{Структура}

\begin{itemize}
    \item in.txt -- SPEC\_FILE, содержащий отрезок и три функции
    \item ast.c и diff\_ast.c -- модули построения Abstract Syntax Tree и его дифференцирования.
    \item rpn2asm.c -- строит дерево функций из SPEC\_FILE, затем на основе дерева и дифференцированного дерева составляет ассемблерные функции в соответствии с cdecl.
    \item funcs.asm -- ассемблерный модуль, составленный rpn2asm.c.
    \item numeric.c -- модуль с реализацией численных методов нахождения корней и интегрирования.
    \item main.c -- ввод и вычисления с применением функций из numeric.c.
    \item Makefile -- сценарий автоматической сборки программы и отчета.
    \item report.tex -- TeX файл с отчетом.
\end{itemize}

\begin{figure}[h]
\centering
\includegraphics[width=15cm]{report/graph.png}
\caption{Графическое представление связей компонентов}
\end{figure}

\subsection{Реализация}
Построение синтаксического дерева само по себе достаточно просто; но, в связи с тем, что в стеке FPU x87 всего 8 регистров, дерево требует оптимизации, чтобы его высота не превышала количество регистров. Оптимизации включают:
\begin{itemize}
    \item Балансировку цепочек ассоциативных операций
    \item Удаление нулевых узлов и единичных множителей (появляется много после дифференцирования)
    \item Сокращение операций вида $\frac{f(x)}{f(x)} := 1$
    \item Предвычисление операций с константами (иногда появляются после дифференцирования)
\end{itemize}
Затем, на основе дерева строятся ассемблерные функции; дерево обходится в post-order порядке, левый операнд лежит в st1, правый в st0; после применения оператора кладем результат в st0.\\

\newpage

\section{Сборка программы (Make-файл)}

Текст Makefile:
\inputminted[frame=single, linenos, breaklines=true, fontsize=\small]{text}{Makefile}

Поскольку изменение флагов (переменных окружения) USE\_X make не считает изменением зависимостей, для удобного изменения используемых методов файлы numeric.o и main.o, зависящие от этих флагов, пересобираются каждый раз. Для этого им задана зависимость force, обновляемая каждый раз.\\\\
Переменные окружения (в скобках -- значение по умолчанию):
\begin{itemize}
    \item USE\_NEWTON=0,1 (0) -- использовать метод Ньютона (касательных) или метод хорд
    \item USE\_SIMPSON=0,1 (1) -- использовать метод Симпсона (парабол) или метод трапеций
    \item SPEC\_FILE=<filename> (in.txt) -- файл с отрезком и кривыми
    \item SAFE\_EXP=0,1 (1) -- использовать упрощенную экспоненту (см. стр. 14)
\end{itemize}

\begin{figure}[h]
\centering
\includegraphics[width=15cm]{report/make_graph.png}
\caption{Графическое представление зависимостей}
\end{figure}





\newpage

\section{Анализ допущенных ошибок}

Задание не включало реализацию оператора \^{}; он был реализован, но с нюансами. В FPU x87 нет прямой инструкции для возведения $a$ в степень $b$; есть инструкция fyl2x, вычисляющая $a \cdot log_2(b)$, и инструкция f2xml, вычисляющая $2^x-1$ для $x \in [-1, 1]$; таким образом, $a^b = 2^{b \cdot log_2(a)} = 2^{round(b \cdot log_2(a))} \cdot 2^{b \cdot log_2(a) \% 1}$ (целая и дробная части).\\

Проблемы возникают, когда $a \leq 0$; $log_2(0) = -\infty$, для отрицательных -- nan. Отрицательный аргумент появляется очень часто, когда вычисляется, например, значение $(x-a)^2$ при $x < a$; поэтому в программе полноценно реализовано возведение только в константную степень, которая превращается в обычное перемножение в ходе построения дерева. В Makefile для включения реальной экспоненты нужно установить SAFE\_EXP=0.

\newpage
\begin{raggedright}
\addcontentsline{toc}{section}{Список цитируемой литературы}
\begin{thebibliography}{99}
\bibitem{math} Ильин~В.\,А., Садовничий~В.\,А., Сендов~Бл.\,Х. Математический анализ. Т.\,1~---
    Москва: Наука, 1985.\\
\bibitem{docs}
        \url{https://artteam8.github.io/1st-Year-Final-Boss}\\
\bibitem{code} 
        \url{https://github.com/artteam8/1st-Year-Final-Boss}\\
\end{thebibliography}
\end{raggedright}

\end{document}
